\subsection{ANTICOMMUTATIVITY OF THE EXTERIOR ALGEBRA}

PROPOSITION 5. (i) Let \((x_{i})_{1\leqslant i\leqslant n}\) be a finite sequence of elements of the module \(\mathbf{M}\) ; for every permutation \(\sigma\) in the symmetric group \(\mathfrak{S}_n\) ,

\[
x _ {\sigma (1)} \wedge x _ {\sigma (2)} \wedge \dots \wedge x _ {\sigma (n)} = \varepsilon_ {\sigma}. x _ {1} \wedge x _ {2} \wedge \dots \wedge x _ {n}\tag{6}
\]

where \(\varepsilon_{\sigma}\) denotes the signature of the permutation \(\sigma\) .

\begin{enumerate}
\def\labelenumi{(\roman{enumi})}
\setcounter{enumi}{1}
\tightlist
\item
  If there exist two distinct indices \(i, j\) such that \(x_{i} = x_{j}\) , the product
\end{enumerate}

\[
x _ {1} \wedge x _ {2} \wedge \dots \wedge x _ {n}
\]

is zero.

\begin{enumerate}
\def\labelenumi{(\roman{enumi})}
\tightlist
\item
  First of all, since \(x \wedge x = 0\) for all \(x \in \mathbf{M}\) by definition of the ideal \(\mathfrak{J}''\) , also, for \(x, y\) in \(\mathbf{M}\) ,
\end{enumerate}

\[
x \wedge y + y \wedge x = (x + y) \wedge (x + y) - x \wedge x - y \wedge y = 0.
\]

This establishes (6) in the case \(n = 2\) . The general case then follows from § 4, no. 6, Lemma 3.

\begin{enumerate}
\def\labelenumi{(\roman{enumi})}
\setcounter{enumi}{1}
\tightlist
\item
  Under the hypothesis of (ii), there exists a permutation \(\sigma \in \mathfrak{S}_n\) such that \(\sigma(1) = i\) and \(\sigma(2) = j\) ; then the left hand side of (6) is zero for this permutation and hence so is the right hand side.
\end{enumerate}

COROLLARY 1. Let H, K be two complementary subsets of the interval \([1, n]\) of N and let \((i_{h})_{1 \leqslant h \leqslant p}, (j_{k})_{1 \leqslant k \leqslant n - p}\) be the sequences of elements of H and K respectively, arranged in increasing order; we write

\[
x _ {H} = x _ {i _ {1}} \wedge x _ {i _ {2}} \wedge \dots \wedge x _ {i _ {p}}, \quad x _ {K} = x _ {j _ {1}} \wedge x _ {j _ {2}} \wedge \dots \wedge x _ {j _ {n - p}};
\]

then

\[
x _ {\mathrm{H}} \wedge x _ {\mathrm{K}} = (- 1) ^ {\nu} x _ {1} \wedge x _ {2} \wedge \dots \wedge x _ {n}\tag{7}
\]

where \(\nu\) is the number of ordered pairs \((i,j)\in \mathbf{H}\times \mathbf{K}\) such that \(i > j\)

By Proposition 5 this reduces to proving

Lemma 1. If \(\sigma \in \mathfrak{S}_n\) is the permutation such that \(\sigma(h) = i_h\) for \(1 \leqslant h \leqslant p\) , \(\sigma(h) = j_{h-p}\) for \(p + 1 \leqslant h \leqslant n\) , then \(\varepsilon_\sigma = (-1)^{\nu}\) .

For \(1 \leqslant h < h' \leqslant p\) or \(p + 1 \leqslant h < h' \leqslant n\) , \(\sigma(h') > \sigma(h)\) and the number of ordered pairs \((h, h')\) such that \(1 \leqslant h \leqslant p < h' \leqslant n\) and \(\sigma(h) > \sigma(h')\) is equal to \(\nu\) .

COROLLARY 2. The graded algebra \(\bigwedge (\mathbf{M})\) is alternating (§ 4, no. 9).

It suffices to apply Proposition 13 of §4, no. 9 to \(\bigwedge (\mathbf{M})\) , taking as generating system the set \(\mathbf{M}\) and using Proposition 5.

PROPOSITION 6. If M is a finitely generated A-module, \(\bigwedge (\mathbf{M})\) is a finitely generated A-module; also, if M admits a generating system with \(n\) elements, then \(\bigwedge^{p}(\mathbf{M}) = \{0\}\) for \(p > n\) .

Let \((x_{i})_{1\leqslant i\leqslant n}\) be a generating system of M. Every element of \(\bigwedge^{p}(\mathbf{M})\) is a linear combination of elements of the form

\[
x _ {i _ {1}} \wedge x _ {i _ {2}} \wedge \dots \wedge x _ {i _ {p}}
\]

where the indices \(i_k\) belong to \([1, n]\) ; by Proposition 5, these indices can be assumed to be distinct (otherwise the corresponding element is zero). If \(p > n\) , there is no such sequence of indices and hence \(\bigwedge^p(\mathbf{M}) = \{0\}\) . If \(p \leqslant n\) , these sequences are finite in number, which completes the proof.

